\documentclass[11pt]{article}

%------------------------- fonts & encoding -----------------------------------
\usepackage[T1]{fontenc}
\usepackage[utf8]{inputenc}
\usepackage{amsmath,amssymb,amsthm}
\usepackage{mathtools}
\usepackage[sc]{mathpazo}
\linespread{1.06}
\usepackage{microtype}

%------------------------- layout ---------------------------------------------
\usepackage[a4paper,margin=1.05in,marginparwidth=1.0in]{geometry}
\usepackage{booktabs}
\usepackage{enumitem}
\setlist{itemsep=2pt,topsep=3pt}
\usepackage{xcolor}
\usepackage{graphicx}
\usepackage{tikz}
\usepackage{pgfplots}
\pgfplotsset{compat=1.17}
\usepackage[most,breakable]{tcolorbox}

%------------------------- palette --------------------------------------------
\definecolor{crimson}{HTML}{A51C30}
\definecolor{ink}{HTML}{1A1A1A}
\definecolor{steel}{HTML}{2C5F8A}
\definecolor{amber}{HTML}{9A6A00}
\definecolor{teal}{HTML}{1F6F6B}
\definecolor{paperlt}{HTML}{F7F4EF}

%------------------------- section styling ------------------------------------
\usepackage{titlesec}
\titleformat{\section}{\Large\bfseries\color{steel}}{\thesection}{0.7em}{}
\titleformat{\subsection}{\large\bfseries\color{ink}}{\thesubsection}{0.6em}{}
\titlespacing*{\section}{0pt}{1.6em}{0.6em}

%------------------------- hyperref -------------------------------------------
\usepackage[colorlinks=true,linkcolor=steel,citecolor=steel,urlcolor=crimson]{hyperref}

%------------------------- theorem environments -------------------------------
\theoremstyle{definition}
\newtheorem{definition}{Definition}[section]
\newtheorem{example}{Example}[section]
\newtheorem{exercise}{Exercise}[section]
\theoremstyle{plain}
\newtheorem{proposition}{Proposition}[section]
\theoremstyle{remark}
\newtheorem*{remark}{Remark}

%------------------------- callout boxes --------------------------------------
\newtcolorbox{intuition}{
  enhanced, breakable, colback=crimson!5, colframe=crimson, boxrule=0.5pt,
  left=8pt,right=8pt,top=6pt,bottom=6pt, arc=2pt,
  colbacktitle=crimson, coltitle=paperlt, fonttitle=\bfseries, title={Intuition}}

\newtcolorbox{pitfall}{
  enhanced, breakable, colback=amber!7, colframe=amber, boxrule=0.5pt,
  left=8pt,right=8pt,top=6pt,bottom=6pt, arc=2pt,
  colbacktitle=amber, coltitle=paperlt, fonttitle=\bfseries, title={Pitfall}}

\newtcolorbox{keyresult}{
  enhanced, breakable, colback=steel!6, colframe=steel, boxrule=0.5pt,
  left=8pt,right=8pt,top=6pt,bottom=6pt, arc=2pt,
  colbacktitle=steel, coltitle=paperlt, fonttitle=\bfseries, title={Takeaway}}

\newtcolorbox{observed}{
  enhanced, breakable, colback=teal!6, colframe=teal, boxrule=0.5pt,
  left=8pt,right=8pt,top=6pt,bottom=6pt, arc=2pt,
  colbacktitle=teal, coltitle=paperlt, fonttitle=\bfseries, title={Measured}}

%------------------------- math operators -------------------------------------
\DeclareMathOperator{\Var}{Var}
\DeclareMathOperator{\Cov}{Cov}
\DeclareMathOperator{\Corr}{corr}
\DeclareMathOperator{\diag}{diag}
\DeclareMathOperator{\std}{std}
\DeclareMathOperator*{\avg}{\mathbb{E}}
\newcommand{\T}{^{\!\top}}
\newcommand{\R}{\mathbb{R}}
\newcommand{\half}{\tfrac12}
\newcommand{\wvec}{\mathbf{w}}
\newcommand{\Xmat}{\mathbf{X}}
\newcommand{\Fmat}{\mathbf{F}}
\newcommand{\Vmat}{\mathbf{V}}
\newcommand{\Umat}{\mathbf{U}}
\newcommand{\Rmat}{\mathbf{R}}
\newcommand{\Lam}{\boldsymbol{\Lambda}}
\newcommand{\Del}{\boldsymbol{\Delta}}
\newcommand{\fvec}{\mathbf{f}}

%------------------------- title ----------------------------------------------
\title{\vspace{-2.2em}
  {\color{steel}\rule{\textwidth}{1.3pt}}\\[0.7em]
  \Huge\bfseries Industry Factors\\[0.35em]
  \large\mdseries The membership layer --- one factor per industry, with exposure by belonging\\[0.4em]
  {\color{steel}\rule{\textwidth}{0.7pt}}}
\author{}
\date{}

%------------------------- watermark footer -----------------------------------
\usepackage{fancyhdr}
\pagestyle{fancy}
\fancyhf{}
\renewcommand{\headrulewidth}{0pt}
\renewcommand{\footrulewidth}{0pt}
\fancyfoot[L]{\footnotesize\color{black!45}MSCI Barra USE4 Imitation}
\fancyfoot[R]{\footnotesize\color{black!45}\thepage}

\begin{document}
\maketitle
\thispagestyle{empty}

\begin{abstract}
The industry factors are the membership layer of the model: one factor for each industry, with a stock's exposure being simply whether it belongs to that industry. Where the style factors are continuous, standardised scores, industry exposures are binary -- a stock carries a $1$ in its own industry's column and a $0$ in every other. The economic content is the most intuitive in the model: firms in the same industry rise and fall together, sharing exposure to the same input costs, customers, regulation, and cycle, and the industry factor's return is exactly that common move. USE4 specifies around sixty industry factors built on a licensed GICS classification, with diversified firms split fractionally across industries using business-segment data. Our build has neither GICS nor segment data, so it groups Sharadar's raw industry labels into a smaller scheme of \textbf{55 factors} and gives every stock a single, whole membership. Two features set this layer apart from the styles: the exposures are dummies, so there is no trimming or standardisation to do, and the factors carry an identification constraint -- their cap-weighted returns sum to zero -- that ties them to the country factor and is essential to making the whole model estimable.
\end{abstract}

\begin{tcolorbox}[colback=paperlt,colframe=ink!30,boxrule=0.4pt,arc=2pt,left=10pt,right=10pt,top=8pt,bottom=8pt]
\textbf{How to read these notes.} We move from \emph{why} industry membership is a factor to \emph{how} USE4 specifies the layer, then to how we build it without a licensed classification, with a deliberate stop at the identification constraint (Section~\ref{sec:constraint}) that makes the industries and the country factor estimable together. Read that section slowly; it carries the weight of the set. The numbers labelled as measured come from a single run of our personal build.

\smallskip
Four colored boxes reoccur: \textcolor{crimson}{\textbf{Intuition}} gives the mental picture, \textcolor{amber}{\textbf{Pitfall}} flags what breaks without the fix, \textcolor{steel}{\textbf{Takeaway}} states a load-bearing result, and \textcolor{teal}{\textbf{Measured}} reports real numbers from the build run.

\smallskip
\textbf{Reference material.}
\begin{itemize}[leftmargin=1.4em,topsep=2pt]
  \item External: Menchero, Orr \& Wang (2011), USE4 Empirical Notes (Appendix A) and Methodology Notes (Sec. 2.1--2.3).
\end{itemize}
\end{tcolorbox}

\tableofcontents
\bigskip

\section{What industry factors are, and why they are factors}\label{sec:what}

\paragraph{Learning objectives.}
\begin{itemize}[leftmargin=1.4em,topsep=2pt]
  \item State industry exposure as binary membership, in contrast to a continuous style score.
  \item Explain why shared industry membership is a source of common return.
  \item Say in one sentence what an industry factor's return represents.
\end{itemize}

Industry membership is the most concrete source of common movement in the market. Two banks face the same yield curve and the same credit cycle; two oil producers ride the same crude price; two chipmakers share a supply chain and a demand wave. Whatever happens to an industry happens, to a first approximation, to all of its members at once. That shared fate is a risk a model must capture before any subtler style tilt, and the industry factors capture it directly: one factor per industry, whose return is the common move shared by the stocks in it.

What makes this layer different from the styles is the shape of the exposure. A style exposure is a continuous, standardised number -- how cheap, how volatile, how large a stock is, measured on a common scale. An industry exposure is a plain yes-or-no: a stock either belongs to an industry or it does not. So the exposure is binary -- a $1$ in the stock's own industry and a $0$ everywhere else -- and there is nothing to centre, scale, or trim. Membership is the measurement.

\begin{intuition}
Think of the industry layer as a set of club memberships. Every stock holds a membership card for exactly one industry; that card is its exposure, worth $1$ for its own club and $0$ for all the others. When the regression later asks ``what did each industry do this month?'', it is simply reading the average return common to each club's members. The styles ask \emph{how much} of a trait a stock has; the industries ask only \emph{which club} it is in.
\end{intuition}

\begin{keyresult}
The industry layer is a block of binary membership indicators in the exposure matrix: each stock loads $1$ on its own industry and $0$ on the rest. There are no continuous scores here, so no standardisation -- the entire design is the classification.
\end{keyresult}

\section*{Notation}
\addcontentsline{toc}{section}{Notation}

\begin{center}
\begin{tabular}{@{}ll@{}}
\toprule
Symbol & Meaning \\
\midrule
$J$ & number of industry factors (here $J=55$) \\
$X^{\text{ind}}_{ij}$ & industry exposure of stock $i$ to industry $j$: $1$ if $i$ belongs to $j$, else $0$ \\
$f^{\text{ind}}_j$ & return of industry factor $j$ (recovered later, in the cross-sectional regression) \\
$w_j$ & cap weight of industry $j$ --- its share of total in-ESTU market cap \\
$\sum_j w_j f^{\text{ind}}_j = 0$ & the identification constraint on the industry returns \\
ESTU & estimation universe (MSCI USA IMI proxy) \\
\bottomrule
\end{tabular}
\end{center}

\section{How USE4 defines the layer}\label{sec:def}

\paragraph{Learning objectives.}
\begin{itemize}[leftmargin=1.4em,topsep=2pt]
  \item State the classification basis, the count, and the membership rule USE4 uses.
  \item Name the identification constraint the industry returns obey.
\end{itemize}

USE4 builds its industry layer on the GICS classification, with roughly sixty industry factors (Menchero, Orr \& Wang 2011, Empirical Notes, App. A). Some stocks truly load $1$ on a single industry, but USE4 does not force every firm into one box: a diversified company is spread \emph{fractionally} across several industries in proportion to its business segments, so a conglomerate that earns half its revenue in one industry and half in another carries a $0.5$ in each. In USE4 a substantial share of market cap -- well over half -- is multi-industry in this way.

The industry returns are not left free. They obey an identification constraint: weighted by industry cap share, they sum to zero,
\[
\sum_{j} w_j\, f^{\text{ind}}_j \;=\; 0 .
\]
This works because the cap‑share weights are positive while the factor returns (for industries or any factor) can be positive or negative, and, for industries, their cap‑weighted average is constrained to zero; for style factors, an analogous zero‑mean condition is applied to the standardized factor scores rather than directly to the returns. This constraint is what lets the industries coexist with a separate country factor, and Section~\ref{sec:constraint} is devoted to why. Crucially, the industry layer is a set of dummies, so the USE4 standardisation machinery -- the three-sigma trim, the cap-weighted-mean-zero scaling -- simply does not apply to it. There is no descriptor to standardise; the exposure is the membership.

\section{How we build it without a licensed classification}\label{sec:build}

\paragraph{Learning objectives.}
\begin{itemize}[leftmargin=1.4em,topsep=2pt]
  \item Explain the two data gaps (no GICS, no segments) and how each is handled.
  \item State how many factors result and the membership rule used.
\end{itemize}

Two ingredients USE4 relies on are missing from public data, and each forces a deviation. \textbf{First, GICS is licensed and absent.} What Sharadar does provide is a set of raw industry labels -- on the order of a hundred and fifty granular tags. The build groups those raw labels into a smaller set of industry factors, guided by loose, sensible principles: granular labels that describe the same business are pooled; groups too thin to estimate reliably are merged with their nearest neighbour; a couple of groups that have swollen too large to behave as one factor are split; and the naturally large, coherent industries are kept whole -- the whole scheme aimed at the published USE4 industry list. The result is a scheme of \textbf{55 industry factors}. (The specific label-to-factor map used in the build run is kept proprietary; what matters for the reader is the shape of the result, not the bookkeeping of which tag went where.)

\textbf{Second, business-segment data is absent}, so fractional membership cannot be computed. Every stock therefore receives a \emph{single, whole} membership: it loads $1$ on its one assigned industry and $0$ on the rest. This is the layer's most material simplification. Where USE4 would spread a diversified firm across several industries, our build must place it in one, so a true conglomerate's risk is attributed entirely to a single industry rather than shared across the businesses it actually runs.

\begin{pitfall}
A classification scheme must fail loudly, not quietly. If the vendor introduces a new raw label the scheme has never seen, the safe response is to halt and have a human place it -- not to sweep it into a catch-all, which would silently misclassify a stock and corrupt its industry return. This is distinct from the fallback bucket, which legitimately holds only a stock the vendor leaves \emph{unlabelled} altogether (at most one in-universe name on any date here): a label that exists but is unrecognised is a question for a person, not a default for the code.
\end{pitfall}

\begin{intuition}
The grouping is best understood as drawing $55$ sensible boxes and dropping each stock in one. Too many boxes and some hold only a handful of names, so their ``industry return'' is just noise; too few and a box mixes businesses that have nothing to do with each other. The art is choosing boxes large enough to estimate but pure enough to mean something --- which is exactly the trade-off the merge-thin, split-bloated, keep-large-whole principles manage.
\end{intuition}

\section{The identification constraint}\label{sec:constraint}

\paragraph{Learning objectives.}
\begin{itemize}[leftmargin=1.4em,topsep=2pt]
  \item Explain why the industries and a country factor are collinear without a constraint.
  \item State what the cap-weighted-sum-zero constraint does and why the cap-weight vector is built.
\end{itemize}

Here is the problem the constraint solves. The model also carries a single \emph{country} factor that every stock loads on -- a column of all ones. But every stock also belongs to exactly one industry, so the industry columns of the exposure matrix add up, row by row, to that same column of all ones. The country factor is therefore an exact linear combination of the industries: the two are collinear, and a regression cannot separately identify a country return and a full set of industry returns. Something has to give.

The fix is the constraint $\sum_j w_j f^{\text{ind}}_j = 0$: the cap-weighted average of the industry returns is pinned to zero. With it, the redundancy is resolved cleanly and the pieces acquire their natural meaning. The country factor absorbs the market-wide move -- the return common to everything -- while each industry factor becomes a pure \emph{relative} move, how that industry did against the cap-weighted market rather than in absolute terms. This is why the build produces, alongside the membership panel, a per-date vector of industry cap weights $w_j$ -- each industry's share of in-ESTU market cap, summing to one across the layer. Those weights are the $w$ in the constraint, and the regression that recovers the factor returns imposes the constraint using them.

\begin{intuition}
Without the constraint, the model cannot tell ``the whole market rose $1\%$'' from ``every industry rose $1\%$'' -- they are the same numbers wearing different labels. Forcing the cap-weighted industry returns to sum to zero declares, by convention, that the market-wide part lives in the country factor and the industries only describe deviations from it. It is a choice of bookkeeping, but a necessary one: it gives each factor a job that does not overlap its neighbour's.
\end{intuition}

\begin{keyresult}
The industry factors are collinear with the country factor because every stock is in some industry. The cap-weighted-sum-zero constraint breaks the tie: the country factor carries the market, the industries carry only relative moves. The cap-weight vector the build stores is precisely the weighting that constraint uses.
\end{keyresult}

\section{The shape of the layer}\label{sec:shape}

\paragraph{Learning objectives.}
\begin{itemize}[leftmargin=1.4em,topsep=2pt]
  \item Describe the membership panel and why no standardisation diagnostics appear.
  \item Read the divergence between an industry's cap weight and its head count.
\end{itemize}

Because the exposures are dummies, the diagnostics that dominate a style factor -- clip rates, skew, kurtosis, the fraction within three standard deviations -- have no meaning here. What characterises the layer instead is its coverage and its balance: every one of the $55$ factors is populated on every date, essentially every stock is classified (the catch-all fallback holds at most one in-universe name on a date), and no single industry is allowed to dominate -- the largest industry's share of market cap stays around one-seventh.

The one genuinely informative cross-sectional picture is how an industry's \emph{cap weight} relates to its \emph{head count}, and the two diverge sharply (Figure~\ref{fig:industries_dist}). A handful of factors carry a large share of market capitalisation with relatively few member names -- industries dominated by a small number of very large firms. Others are populous, holding many names, yet command only a modest share of cap because their members are individually small. The cap weight, not the head count, is what enters the identification constraint, so a factor with few names can still pull hard on the model.

\begin{observed}
On the build run the industry layer is a $0/1$ membership panel of $\approx 1.77$M stock-month rows over $330$ dates and $\approx 16{,}800$ tickers, with $55$ factors present on every date and at most one in-ESTU stock unclassified on any date. There is no trim and no standardisation -- the exposures are dummies. The per-date cap-weight vector sums to one across the $55$ factors, and the single largest industry never exceeds $\approx 14\%$ of in-ESTU cap. Cap weight and head count are only loosely related (Figure~\ref{fig:industries_dist}): the most cap-heavy factors are not the most populous, and vice versa.
\end{observed}

\begin{figure}[htbp]
\centering
\includegraphics[width=\linewidth]{figures/industries_dist.png}
\caption{Each point is one of the $55$ industry factors on a single date: its in-ESTU head count (horizontal) against its share of in-ESTU market cap (vertical). The two diverge -- a few factors carry a large share of cap with relatively few names (top-left), while the most populous factors carry only modest cap (lower-right). The largest single industry sits near $14\%$ of cap; most factors are well below $4\%$. Cap weight, the quantity that enters the identification constraint, is not headcount.}
\label{fig:industries_dist}
\end{figure}

\section{Summary, and one caveat}\label{sec:summary}

\paragraph{Learning objectives.}
\begin{itemize}[leftmargin=1.4em,topsep=2pt]
  \item Recite how the industry layer is built and constrained.
  \item State the central limitation and what would address it.
\end{itemize}

The build, end to end: group Sharadar's raw industry labels into a scheme of $55$ industry factors; assign every stock a single, whole membership, so its industry exposure is a $1$ in one column and $0$ in the rest; build the per-date cap-weight vector that the identification constraint will use; and hand the membership panel and the weights downstream, where the cross-sectional regression recovers the industry returns subject to $\sum_j w_j f^{\text{ind}}_j = 0$. There is no trimming or standardising, because the exposures are dummies, not descriptors.

\begin{keyresult}
Industries are the model's membership layer: binary exposures, no standardisation, and an identification constraint that lets them sit alongside the country factor. The two computations that matter are the classification itself and the cap-weight vector the constraint runs on.
\end{keyresult}

\begin{pitfall}
One honest limitation dominates: single membership. Without business-segment data the build cannot spread a diversified firm across the industries it actually operates in, so a true conglomerate's risk is loaded entirely onto one industry -- exactly the firms USE4 treats fractionally, and exactly where a single label is least faithful. Two smaller points compound it: the classification is fewer factors than USE4's (a coarser partition, from grouping a limited set of public labels), and it is applied as a current static snapshot rather than point-in-time, so a firm that has been reclassified over the years carries its present label backward through history. The clear levers are a licensed GICS classification, business-segment data to restore fractional membership, and a point-in-time classification history.
\end{pitfall}

\begin{remark}
None of this changes the layer's place in the model: the industries are the membership backbone of the exposure matrix, capturing the most concrete form of common risk and, through their identification constraint, making the country factor and the cross-sectional regression well posed. The caveats are about \emph{resolution} -- how finely the classification cuts and whether it can split a firm -- not about whether industry membership is a real and first-order risk.
\end{remark}

\begin{center}\small\itshape
References: Menchero, Orr \& Wang (2011), \textit{The Barra US Equity Model (USE4)} --- Empirical Notes (Appendix A) and Methodology Notes (Sec. 2.1--2.3). ESTU treated as an MSCI USA IMI proxy; measured quantities are from a personal build run.
\end{center}

\end{document}
